\documentclass[11pt,a4paper]{article}
\usepackage[margin=27mm]{geometry}
\usepackage{amsmath,amssymb,amsthm,hyperref}
\emergencystretch=1em
\newtheorem{theorem}{Theorem}\newcommand{\Sz}{\operatorname{Sz}}
\newcommand{\F}{\mathbb F}
\title{Scalar subspaces and flag transitivity}
\author{Human Compression project}\date{}
\begin{document}\maketitle
\begin{abstract}
The previously acquired scalar-subspace family satisfies the literal
algebraic local two-transitivity condition in arbitrary rank. Combining
its intersection formula with the proved descent theorem shows that
these pregeometries are thick, residually connected, flag-transitive
geometries. The additional argument is a surjectivity calculation on
independent scalar coordinates.
\end{abstract}

Let $e\ge3$ be odd, put $t=2^e$ and $K=\F_t$, and choose $k\ge2$
distinct odd primes $a_1,\ldots,a_k$. Set
\[
 A=\prod_i a_i,\quad F=\F_{t^A},\quad
 F_i=\F_{t^{A/a_i}},\quad G=\Sz(t^A),\quad S_i=\Sz(F_i).
\]
Choose $x_i\in\F_{t^{a_i}}\setminus K$ and put
\[
 T=\operatorname{span}_K(1,x_1,\ldots,x_k),\qquad
 Q=Z_TB_t,
\]
where $B_t=U_tH_t$ is the standard point stabilizer of $\Sz(t)$ and
$Z_T$ is the central unipotent subgroup with parameter in $T$.
The earlier scalar-subspace proof establishes that $1,x_1,\ldots,x_k$
are independent over $K$ and, for $J\subseteq\{1,\ldots,k\}$,
\[
 T\cap F_J=T_J=\operatorname{span}_K(1,x_i:i\notin J),\qquad
 Q\cap\bigcap_{j\in J}S_j=Z_{T_J}B_t,
 \tag{1}
\]
where $F_J=\F_{t^{A/\prod_{j\in J}a_j}}$.
For completeness, apply $\sigma_i:y\mapsto y^{t^{A/a_i}}$ to a
linear combination lying in $F_i$. This automorphism fixes every $x_j$
with $j\ne i$, fixes $K$, and moves $x_i$. Subtracting the two
expressions forces its $x_i$ coefficient to vanish. The same argument
with a zero linear combination proves independence, and intersection
over $i\in J$ gives the displayed formula.
The coordinate multiplication law makes $Z_TB_t$ a subgroup: $U_t$
centralizes $Z_T$, and $H_t$ acts by all nonzero $K$-scalars.

\begin{theorem}
The coset system $\Gamma(G;S_1,\ldots,S_k,Q)$ is a thick,
residually connected, flag-transitive geometry of rank $k+1$.
It satisfies the source's residual weak primitivity condition and
local two-transitivity, now on its actual rank-one residues.
\end{theorem}
\begin{proof}
Put $T_i=\operatorname{span}_K(1,x_j:j\ne i)$ and
$H_i=Q\cap S_i=Z_{T_i}B_t$. Every left $H_i$-coset in $Q$ has a
representative in $Z_T$, and
\[
 T/T_i\longrightarrow Q/H_i,\qquad b+T_i\longmapsto z(b)H_i
 \tag{2}
\]
is a bijection of sets. Indeed $Q=Z_TH_i$ and
$Z_T\cap H_i=Z_{T_i}$, by (1). No quotient-group structure on
$Q/H_i$ is needed. Given arbitrary coordinates $c_i\in K$, the
single vector $\sum_i c_ix_i\in T$ represents them in every $T/T_i$.
Thus
\[
 Q\longrightarrow\prod_i Q/H_i
\]
is surjective.

The standard subfield cosets have the common-representative property by
the acquired descent theorem, with Steinberg endomorphism $\rho^e$
and independent relative degrees $a_i$. Its attachment lemma and (2)
therefore give that property for all the types including $Q$.
Every flag extends to a translate of the standard chamber; hence this
is a flag-transitive geometry. In particular all its actual residues
are isomorphic to the standard subgroup-intersection residues.

The earlier proof of residual weak primitivity can now be read in these
actual residues. Before selecting $Q$, an unused field type makes a
maximal prime-degree inclusion of simple Suzuki groups, since the bottom
field exponent is at least $e>1$. After every field type has been
selected, the remaining inclusion is $B_t<\Sz(t)$, which is maximal.
Once $Q$ is selected, an unused field type replaces $T_J$ by a
hyperplane containing $K$. The overgroups of $Z_{T_J\cap F_i}B_t$
inside $Z_{T_J}B_t$ correspond to $K$-subspaces between those two
parameter spaces: the proof is the elementary equality
$L=(L\cap Z_{T_J})B_t$, together with the scalar action of $H_t$.
The codimension is one, so every such inclusion is maximal.

The full intersection is $B_t$. The panel missing $Q$ has
$[\Sz(t):B_t]=t^2+1$ chambers and the usual doubly transitive Suzuki
action. The panel missing $S_i$ has $[Z_{K+Kx_i}B_t:B_t]=t$
chambers. Its quotient action on $(K+Kx_i)/K$ contains all translations
and all nonzero $K$-scalars; it is the doubly transitive affine action
on $K$. Since $t\ge8$, every panel has at least three chambers.
This proves thickness and actual local two-transitivity.

Finally consider a standard residue of rank at least two. If $Q$ has
been selected, at least two field types remain. Their two parameter
hyperplanes sum to $T_J$, so their subgroups generate $Z_{T_J}B_t$.
If $Q$ is unselected and at least two field types remain, two distinct
maximal simple subfield groups generate the current Suzuki group.
If precisely one field type and $Q$ remain, the base $\Sz(t)$ is
maximal in $\Sz(t^{a_i})$, and $Z_{K+Kx_i}B_t$ is not contained in
it because $x_i\notin K$. Again they generate the residue group.
The coset incidence graph is connected exactly when its parabolics
generate; flag transitivity transports this conclusion to every residue.
Thus the geometry is residually connected.
\end{proof}

\paragraph{Attribution and source scope.}
The scalar family and its intersection proof were already independently
accepted in the project; the Frobenius-isolation step there was suggested
by the root reviewer. The descent and coset-attachment lemmas are earlier
project proofs retaining classical Lang--Steinberg. The new step here is
the explicit coordinate bijection (2) and its geometric consequences.
It uses no new finite acquisition and does not change any frozen proof.

The fixed higher-rank source, D.~Leemans, Note di Matematica 20(1)
(2000/2001), p.~2, quotes the earlier assertion that no rank at least four
geometry satisfies both residual weak primitivity and local
two-transitivity. Taking $k\ge3$ disproves this statement even with
firmness, residual connectedness and flag transitivity imposed. For
$k=2$, the rank is three while the ambient field degree is composite,
and every panel is thick. These are further consequences of the same
already acquired family, not new original or principal-refutation counts.
No full Lean formalization or complete classification is claimed.
\end{document}
